\documentclass[10pt,letterpaper]{article}

% ---------- Document setup ----------
\usepackage[letterpaper,top=0.62in,bottom=0.62in,left=0.72in,right=0.72in,
            headheight=13.6pt,footskip=24pt]{geometry}
\usepackage[T1]{fontenc}
\usepackage[utf8]{inputenc}
\usepackage{lmodern}
\usepackage{microtype}
\usepackage[dvipsnames]{xcolor}
\usepackage{enumitem}
\usepackage{titlesec}
\usepackage{tabularx}
\usepackage{array}
\usepackage{needspace}
\usepackage{fancyhdr}
\usepackage{lastpage}
\usepackage[hidelinks]{hyperref}

% Make copied PDF text machine-readable and ATS-friendly.
\input{glyphtounicode}
\pdfgentounicode=1

\hypersetup{
  pdftitle={Daniela Ruiz, Ph.D. -- Academic Curriculum Vitae},
  pdfauthor={Daniela Ruiz},
  pdfsubject={Adjunct Professor of Sociology and Applied Social Research},
  pdfkeywords={adjunct professor, sociology, higher education, teaching, curriculum design, applied social research}
}

% ---------- Visual system ----------
\definecolor{Navy}{HTML}{18324A}
\definecolor{Teal}{HTML}{2D7472}
\definecolor{Gold}{HTML}{B88A3B}
\definecolor{Ink}{HTML}{26323A}
\definecolor{Muted}{HTML}{5E6B73}
\definecolor{Rule}{HTML}{CFD8DC}
\definecolor{Pale}{HTML}{F3F6F6}

\color{Ink}
\setlength{\parindent}{0pt}
\setlength{\parskip}{0pt}
\setlength{\tabcolsep}{0pt}
\urlstyle{same}
\raggedbottom

\titleformat{\section}
  {\large\bfseries\color{Navy}}
  {}{0pt}{}
  [\vspace{2pt}\color{Teal}\titlerule]
\titlespacing*{\section}{0pt}{11pt}{6pt}

\titleformat{\subsection}
  {\normalsize\bfseries\color{Teal}}
  {}{0pt}{}
\titlespacing*{\subsection}{0pt}{8pt}{3pt}

\setlist[itemize]{
  leftmargin=1.25em,
  label=\textcolor{Teal}{\small\textbullet},
  itemsep=2.1pt,
  topsep=3pt,
  parsep=0pt,
  partopsep=0pt
}
\setlist[enumerate]{
  leftmargin=2.15em,
  label=\textcolor{Teal}{\arabic*.},
  itemsep=3.2pt,
  topsep=3pt,
  parsep=0pt,
  partopsep=0pt
}

% ---------- Running matter ----------
\pagestyle{fancy}
\fancyhf{}
\renewcommand{\headrulewidth}{0pt}
\renewcommand{\footrulewidth}{0pt}
\fancyhead[L]{\scriptsize\color{Muted}\textsc{Daniela Ruiz, Ph.D.}}
\fancyhead[R]{\scriptsize\color{Muted}\textsc{Academic Curriculum Vitae}}
\fancyfoot[C]{\scriptsize\color{Muted}Page \thepage\ of \pageref*{LastPage}}

% ---------- Reusable CV commands ----------
\newcommand{\contactsep}{\enspace\textcolor{Gold}{\textbullet}\enspace}

\newcommand{\appointment}[4]{%
  \Needspace{5\baselineskip}%
  \textbf{\color{Navy}#1}\hfill\textbf{#2}\\[-1pt]
  \textit{#3}\hfill\textit{\color{Muted}#4}\par
}

\newcommand{\degree}[4]{%
  \Needspace{3\baselineskip}%
  \textbf{\color{Navy}#1}\hfill\textbf{#2}\\[-1pt]
  #3\hfill\textit{\color{Muted}#4}\par
}

\newcommand{\talk}[3]{%
  \Needspace{3\baselineskip}%
  \textbf{#1}\hfill\textbf{#2}\\[-1pt]
  \textit{#3}\par
}

\newcommand{\referenceentry}[3]{%
  \textbf{\color{Navy}#1}\newline
  #2\newline
  \href{mailto:#3}{#3}%
}

\newcolumntype{Y}{>{\raggedright\arraybackslash}X}

\begin{document}

% ---------- Identity ----------
\begin{center}
  {\fontsize{25}{28}\selectfont\bfseries\color{Navy} DANIELA RUIZ, Ph.D.}\par
  \vspace{2pt}
  {\large\bfseries\color{Teal} Adjunct Professor of Sociology \textcolor{Gold}{|} Applied Social Research}\par
  \vspace{5pt}
  Chicago, Illinois
  \contactsep \href{tel:+13125550142}{+1 312 555 0142}
  \contactsep \href{mailto:daniela.ruiz@protonmail.com}{daniela.ruiz@protonmail.com}\\[-1pt]
  \href{https://www.linkedin.com/in/daniela-ruiz-sociology}{linkedin.com/in/daniela-ruiz-sociology}
  \contactsep \href{https://orcid.org/0000-0002-6841-7305}{ORCID 0000-0002-6841-7305}
  \contactsep \href{https://www.danielaruiz.org}{danielaruiz.org}
\end{center}

\vspace{2pt}
\color{Gold}\rule{\textwidth}{1.2pt}
\color{Ink}

\section*{Academic Profile}

Student-centered sociology educator with 11 years of experience teaching undergraduate learners at public universities and community colleges. Instructor of record for 46 course sections across face-to-face, hybrid, and online formats, with sustained teaching evaluations averaging 4.8/5.0. Expertise in introductory sociology, social inequality, race and ethnicity, research methods, and urban communities. Known for clear course design, inclusive discussion practices, transparent assessment, and applied projects that connect sociological evidence to students' work and civic lives. Active scholar-practitioner with peer-reviewed research on first-generation college access and neighborhood change.

\section*{Teaching Areas}

\begin{tabularx}{\textwidth}{@{}>{\bfseries\color{Navy}}p{1.34in}Y@{}}
Core Sociology & Introduction to Sociology; Social Problems; Sociological Theory; Sociology of the Family \\
Inequality & Race and Ethnicity; Social Stratification; Gender and Society; Immigration; Education and Society \\
Methods & Social Research Methods; Qualitative Methods; Program Evaluation; Community-Based Research \\
Applied Topics & Urban Sociology; Work and Organizations; Community Change; Public Sociology \\
Teaching Formats & In-person; synchronous and asynchronous online; hybrid; accelerated 8-week; dual-enrollment
\end{tabularx}

\section*{Education}

\degree{Ph.D., Sociology}{2015}{University of Illinois Chicago, Chicago, IL}{Urban Sociology and Education}
\vspace{3pt}
Dissertation: \textit{Pathways Through the City: Institutions, Neighborhoods, and First-Generation College Persistence}\\[-1pt]
Committee Chair: Teresa M. Alvarez, Ph.D.

\vspace{5pt}
\degree{M.A., Sociology}{2010}{DePaul University, Chicago, IL}{Applied Social Research}
\vspace{3pt}
\degree{B.A., Sociology, \textit{magna cum laude}}{2007}{Loyola University Chicago, Chicago, IL}{Minor: Latino Studies}

\section*{Academic Appointments}

\appointment{Adjunct Professor of Sociology}{2018--Present}{Lakeshore Metropolitan University, Chicago, IL}{Department of Sociology}
\begin{itemize}
  \item Teach 3--4 sections per academic year, including Introduction to Sociology, Race and Ethnicity, Urban Sociology, and Social Research Methods; typical enrollments range from 24 to 62 students.
  \item Redesigned Introduction to Sociology around low-stakes retrieval practice, scaffolded writing, and locally grounded data exercises; D/F/withdrawal rates declined from 21\% to 12\% across six sections.
  \item Built accessible learning modules in Canvas using backward design, captioned media, screen-reader-ready documents, and transparent assignment prompts aligned to Universal Design for Learning principles.
  \item Earned mean student ratings of 4.8/5.0 for instructor effectiveness and 4.9/5.0 for respect for diverse perspectives across the most recent 12 evaluated sections.
\end{itemize}

\appointment{Adjunct Faculty, Social Sciences}{2015--Present}{Harbor City College, Evanston, IL}{Division of Liberal Arts}
\begin{itemize}
  \item Deliver transfer-level sociology courses in 16-, 12-, and 8-week formats for a diverse student population that includes first-generation learners, working adults, and multilingual students.
  \item Serve as lead adjunct for the common Introduction to Sociology course, coordinating shared outcomes, assessment rubrics, textbook review, and onboarding for seven part-time instructors.
  \item Introduced an open educational resource course pack that eliminated approximately \$28,000 in annual student textbook costs while maintaining department-wide learning-outcome performance.
  \item Piloted an early-alert and structured check-in process that increased assignment completion by 14 percentage points in asynchronous online sections.
\end{itemize}

\appointment{Visiting Lecturer}{2016--2018}{Northbridge University, Chicago, IL}{College of Professional Studies}
\begin{itemize}
  \item Taught evening and weekend sections of Social Problems, Work and Organizations, and Program Evaluation for adult degree-completion students.
  \item Partnered with local nonprofits to create client-based evaluation projects; student teams delivered logic models, survey instruments, and findings briefs to six community organizations.
\end{itemize}

\appointment{Graduate Instructor and Teaching Assistant}{2009--2015}{University of Illinois Chicago, Chicago, IL}{Department of Sociology}
\begin{itemize}
  \item Served as instructor of record for four sections of Introduction to Sociology and as teaching assistant for Sociological Theory, Statistics for the Social Sciences, and Urban Sociology.
  \item Led weekly discussion and methods labs, advised students on research design and academic writing, and received the department's Outstanding Graduate Instructor Award in 2014.
\end{itemize}

\section*{Selected Teaching Effectiveness}

\begin{tabularx}{\textwidth}{@{}>{\bfseries\color{Navy}}p{2.02in}Y>{\raggedleft\arraybackslash}p{0.94in}@{}}
Instructor effectiveness & Mean across 12 recent sections (institutional student survey) & 4.8/5.0 \\
Inclusive learning climate & ``Respects perspectives different from their own'' & 4.9/5.0 \\
Learning-outcome attainment & Students meeting or exceeding common-course benchmark & 86\% \\
Online course quality & Internal Quality Matters-aligned review, SOC 101 & Met all 42 criteria \\
Peer observation & ``Purposeful, well-paced, and notably effective at drawing quieter students into analysis'' & Excellent
\end{tabularx}

\section*{Selected Course Portfolio}

\begin{tabularx}{\textwidth}{@{}>{\bfseries\color{Navy}}p{2.05in}Y>{\raggedleft\arraybackslash}p{0.78in}@{}}
Introduction to Sociology & OER redesign; data labs using census indicators; cumulative social-structure portfolio & 22 sections \\
Race and Ethnicity & Comparative case studies; media discourse analysis; community oral-history assignment & 8 sections \\
Social Research Methods & Mixed-methods proposal; survey design; interviewing; ethics; introductory analysis in R & 6 sections \\
Urban Sociology & Neighborhood observation fieldnotes; housing-policy simulation; Chicago community profile & 5 sections \\
Social Problems & Evidence briefs; structured deliberation; public-facing final project & 3 sections \\
Program Evaluation & Community partner project; logic models; measurement plans; executive presentation & 2 sections
\end{tabularx}

\Needspace{14\baselineskip}
\section*{Curriculum and Pedagogical Leadership}

\begin{itemize}
  \item Co-led departmental alignment of 28 course sections to five measurable general-education outcomes; developed common analytic-writing rubric and facilitated two calibration workshops for full- and part-time faculty.
  \item Created a reusable bank of 40 inclusive, discussion-ready case studies with facilitation notes and alternative asynchronous activities; materials are used across nine sociology and social-science courses.
  \item Designed a four-stage research assignment sequence that reduced last-week submissions and improved average rubric performance in evidence use from 3.1 to 3.8 on a 4-point scale.
  \item Mentor new adjunct instructors on pacing, academic integrity, accessible course materials, difficult-dialogue facilitation, and equitable grading; supported nine colleagues since 2020.
  \item Facilitate annual workshops on transparent assignment design and practical strategies for belonging in commuter and online classrooms.
\end{itemize}

\section*{Scholarship}

\subsection*{Peer-Reviewed Publications}
\begin{enumerate}
  \item \textbf{Ruiz, D.} (2024). Institutional navigators and the uneven work of belonging among first-generation commuter students. \textit{Journal of College Student Development}, 65(4), 421--439.
  \item Patel, N., \textbf{Ruiz, D.}, and Okafor, M. (2022). Teaching the city with the city: Community-engaged inquiry in introductory sociology. \textit{Teaching Sociology}, 50(3), 244--258.
  \item \textbf{Ruiz, D.} and Chen, L. (2020). Beyond displacement: Social infrastructure and resident adaptation in changing neighborhoods. \textit{City \& Community}, 19(4), 982--1005.
  \item \textbf{Ruiz, D.} (2018). Navigating transfer: Organizational signals and first-generation students' institutional choices. \textit{Community College Review}, 46(2), 176--198.
  \item Alvarez, T. M., \textbf{Ruiz, D.}, and Williams, J. (2016). Neighborhood context and college persistence in a metropolitan public university. \textit{Sociology of Education}, 89(3), 214--233.
\end{enumerate}

\subsection*{Public Scholarship and Instructional Resources}
\begin{enumerate}
  \item \textbf{Ruiz, D.} (2023). \textit{Reading Neighborhood Change: An Open Data Workbook for Undergraduate Sociology}. Lakeshore Open Press. CC BY 4.0.
  \item \textbf{Ruiz, D.} (2021). Making introductory research methods consequential through community questions. \textit{The Society Pages: Teaching \& Learning}.
  \item \textbf{Ruiz, D.} (2019). \textit{Chicago in Social Context}, an open collection of 18 teaching cases. Harbor City College OER Repository.
\end{enumerate}

\section*{Selected Presentations and Invited Workshops}

\talk{``Small Changes, Stronger Persistence: Course Design for Commuter Students''}{2025}{Invited workshop, Lakeshore Metropolitan University Center for Teaching Excellence, Chicago, IL}
\vspace{4pt}
\talk{``Community Questions as the Spine of an Introductory Methods Course''}{2024}{Annual Meeting of the American Sociological Association, Montréal, QC}
\vspace{4pt}
\talk{``OER Beyond Cost Savings: Building a More Local Introductory Sociology Course''}{2023}{Midwest Sociological Society Annual Meeting, Minneapolis, MN}
\vspace{4pt}
\talk{``Designing Equitable Participation in Hybrid Discussion Courses''}{2022}{Illinois Online Learning Summit, Virtual}
\vspace{4pt}
\talk{``Institutional Navigators and First-Generation Student Belonging''}{2021}{Association for the Study of Higher Education Annual Conference, San Juan, PR}

\section*{Academic and Professional Service}

\begin{tabularx}{\textwidth}{@{}Y>{\raggedleft\arraybackslash}p{1.10in}@{}}
Lead Adjunct, Sociology Curriculum Committee, Harbor City College & 2021--Present \\
Faculty Advisor, Students for Community Research, Lakeshore Metropolitan University & 2020--Present \\
Member, General Education Assessment Working Group, Harbor City College & 2019--2024 \\
Reviewer, \textit{Teaching Sociology} and \textit{Journal of Applied Social Science} & 2020--Present \\
Volunteer Research Advisor, Chicago Neighborhood Learning Collaborative & 2018--Present \\
Session Organizer, Midwest Sociological Society Annual Meeting & 2022, 2024
\end{tabularx}

\Needspace{9\baselineskip}
\section*{Awards and Professional Development}

\begin{tabularx}{\textwidth}{@{}Y>{\raggedleft\arraybackslash}p{0.72in}@{}}
Excellence in Part-Time Faculty Teaching Award, Harbor City College & 2024 \\
Inclusive Teaching Fellowship, Lakeshore Metropolitan University & 2022 \\
Open Education Innovation Grant, Illinois Community College Consortium, \$8,500 & 2021 \\
Outstanding Graduate Instructor Award, University of Illinois Chicago & 2014 \\
Quality Matters, Applying the Quality Matters Rubric certification & 2023 \\
Association of College and University Educators, Effective College Instruction certificate & 2020
\end{tabularx}

\section*{Research, Instructional, and Language Skills}

\begin{tabularx}{\textwidth}{@{}>{\bfseries\color{Navy}}p{1.34in}Y@{}}
Learning Platforms & Canvas, Blackboard, D2L Brightspace, Moodle, Zoom, Hypothesis, Perusall, Turnitin \\
Research Methods & Interviews, focus groups, survey design, program evaluation, community-based participatory research, IRB protocols \\
Analysis & R, SPSS, Stata, Qualtrics, NVivo, Tableau, U.S. Census and IPUMS data \\
Course Design & Backward design, Universal Design for Learning, transparent assignment design, OER curation, authentic assessment \\
Languages & English and Spanish (bilingual professional proficiency)
\end{tabularx}

\section*{Professional Memberships}

American Sociological Association \contactsep Midwest Sociological Society \contactsep Sociologists for Women in Society\\
National Institute for Staff and Organizational Development \contactsep POD Network in Higher Education

\section*{References}

\begin{tabularx}{\textwidth}{@{}Y Y Y@{}}
\referenceentry{Marisol Chen, Ph.D.}{Chair, Department of Sociology\newline Lakeshore Metropolitan University}{marisol.chen@lakeshoremetro.edu}
&
\referenceentry{Peter Okafor, Ph.D.}{Dean, Division of Liberal Arts\newline Harbor City College}{peter.okafor@harborcity.edu}
&
\referenceentry{Teresa M. Alvarez, Ph.D.}{Professor of Sociology\newline University of Illinois Chicago}{talvarez@uic.edu}
\end{tabularx}

\vfill
\begin{center}
  {\scriptsize\color{Muted}Updated July 2026}
\end{center}

\end{document}
